\section*{Capítulo \ref{ch:b2:enolates-aldol} --- Enoles, enolatos y reacción aldólica}

\begin{solution}{exo:b2:enolates-aldol:1}
Butanona: but-1-en-2-ol \ce{CH2=C(OH)CH2CH3} y but-2-en-2-ol \ce{CH3C(OH)=CHCH3} ((\textit{E}) y
(\textit{Z})). Ciclohexanona: ciclohex-1-en-1-ol.
\end{solution}

\begin{solution}{exo:b2:enolates-aldol:2}
Etano $<$ propanona $<$ 3-oxobutanoato de etilo (10.7) $<$ pentano-2,4-diona (9.0).
\end{solution}

\begin{solution}{exo:b2:enolates-aldol:3}
3-Hidroxibutanal, y luego but-2-enal al calentar.
\end{solution}

\begin{solution}{exo:b2:enolates-aldol:4}
Propanodioato de dietilo, etóxido de sodio, 1-bromobutano; hidrólisis al ácido butilpropanodioico;
el calentamiento elimina \ce{CO2}: \ce{CH3(CH2)3CH2COOH}, el ácido hexanoico.
\end{solution}

\begin{solution}{exo:b2:enolates-aldol:5}
Cinético: desprotonación en C6 (el lado del \ce{CH2}), con LDA a \qty{-78}{\degreeCelsius}. Termodinámico:
el doble enlace más sustituido C1=C2, hacia el carbono que lleva el metilo, con un alcóxido en su
alcohol a temperatura ambiente.
\end{solution}

\begin{solution}{exo:b2:enolates-aldol:6}
4-Fenilbut-3-en-2-ona (bencilidenacetona), \ce{C6H5CH=CHCOCH3}. El benzaldehído no tiene ningún hidrógeno $\alpha$:
no puede formar un \omterm{def:b2:enolates-aldol:enolate}{enolato} y solo es el electrófilo.
\end{solution}

\begin{solution}{exo:b2:enolates-aldol:7}
El \omterm{def:b2:enolates-aldol:enolate}{enolato} de uno de los grupos metilo (C1) ataca el otro carbonilo (C5): un anillo de cinco eslabones,
la 3-metilciclopent-2-en-1-ona tras la deshidratación. La alternativa, el \omterm{def:b2:enolates-aldol:enolate}{enolato} en C3 atacando C5,
formaría un anillo tensionado de tres eslabones.
\end{solution}

\begin{solution}{exo:b2:enolates-aldol:8}
\ce{2CH3COOC2H5 -> CH3COCH2COOC2H5 + C2H5OH} con etóxido de sodio; el producto, desprotonado, se
recupera como 3-oxobutanoato de etilo al acidular.
\end{solution}

\begin{solution}{exo:b2:enolates-aldol:9}
Propanona, etanal y acetofenona (tienen un grupo \ce{CH3-CO}; el etanal es \ce{CH3CHO}).
La pentan-3-ona no.
\end{solution}

\begin{solution}{exo:b2:enolates-aldol:10}
El \omterm{def:b2:enolates-aldol:enolate}{enolato} de un extremo ataca el éster del otro: 2-oxociclopentano-1-carboxilato de etilo, un
anillo de cinco eslabones.
\end{solution}

\begin{solution}{exo:b2:enolates-aldol:11}
En la 3-metilpentan-2-ona el estereocentro (C3) es el carbono $\alpha$: la enolización lo hace plano, y
la reprotonación por una u otra cara lo racemiza. En la 4-metilhexan-2-ona el estereocentro está en C4,
en $\beta$ respecto al carbonilo, y la enolización no lo afecta.
\end{solution}

\begin{solution}{exo:b2:enolates-aldol:12}
Propanona con dos equivalentes de benzaldehído en medio básico. Con dos aldehídos distintos, condénsese primero la
propanona con uno de ellos (un equivalente, aislando la bencilidenacetona), y condénsese luego su
grupo metilo restante con el segundo. En un solo recipiente los dos aldehídos compiten y dan una mezcla de
tres dienonas.
\end{solution}

\begin{solution}{pb:b2:enolates-aldol:1}
\textbf{1.} La propanona, seis hidrógenos $\alpha$ (tres en cada metilo).
\textbf{2.} $\mathrm{CH_3{-}CO{-}CH_2^-} \leftrightarrow \mathrm{CH_3{-}C(O^-){=}CH_2}$.
\textbf{3.} No tiene ningún hidrógeno en el carbono contiguo a su carbonilo (ese carbono pertenece al anillo y
no lleva H).
\textbf{4.} El benzaldehído, un aldehído en exceso, es el mejor electrófilo; el \omterm{def:b2:enolates-aldol:enolate}{enolato} de la propanona lo encuentra
más a menudo que un carbonilo de propanona.
\textbf{5.} Una cetona tiene dos grupos carbonados que ceden densidad electrónica a su carbono carbonílico y
dificultan el acceso; un aldehído tiene uno.
\textbf{6.} Solo hace falta una pequeña fracción de \omterm{def:b2:enolates-aldol:enolate}{enolato}: se regenera a medida que reacciona.
\textbf{7.} El carbono del \omterm{def:b2:enolates-aldol:enolate}{enolato} se une al carbono carbonílico del benzaldehído: el alcóxido de la
4-hidroxi-4-fenilbutan-2-ona.
\textbf{8.} Arranque de un hidrógeno $\alpha$ y pérdida de hidróxido del carbono $\beta$: 4-fenilbut-3-en-2-ona.
\textbf{9.} El doble enlace formado está conjugado a la vez con el anillo y con el carbonilo.
\textbf{10.} El otro grupo metilo aún tiene tres hidrógenos $\alpha$.
\textbf{11.} \ce{2C6H5CHO + CH3COCH3 -> C6H5CH=CHCOCH=CHC6H5 + 2H2O}.
\textbf{12.} Las deshidrataciones, y la precipitación del producto, que abandona la disolución.
\textbf{13.} Dos, ambos (\textit{E}).
\textbf{14.} Los isómeros (\textit{Z}) sitúan el grupo fenilo junto al grupo carbonilo: tensión estérica.
\textbf{15.} Dos \ce{C=C}, un \ce{C=O} y los dos anillos: un sistema conjugado que recorre toda la molécula.
\textbf{16.} Su conjugación extendida estrecha la separación $\pi$ (\cref{prop:b2:huckel:gap}): absorbe en el
extremo azul del visible y se ve amarillo.
\textbf{17.} No: es plano (o casi) y tiene un plano de simetría.
\textbf{18.} 106.12, 58.08 y \qty{234.30}{g/mol}.
\textbf{19.} Benzaldehído: $2.1 \times 1.045/106.12 = \qty{20.7}{mmol}$; propanona: $0.75 \times
0.7845/58.08 = \qty{10.1}{mmol}$.
\textbf{20.} La propanona necesita \qty{20.3}{mmol} de benzaldehído, y se dispone de \qty{20.7}{mmol}:
la propanona es el limitante, por poco.
\textbf{21.} $0.01013 \times 234.30 = \qty{2.37}{g}$.
\textbf{22.} El benzaldehído sobrante queda en el etanol y se elimina con los lavados.
\textbf{23.} El producto de monocondensación, la bencilidenacetona.
\textbf{24.} $0.75 \times 2.37 = \boldsymbol{\approx \qty{1.78}{g}}$ de dibencilidenacetona.
\end{solution}
